\documentclass[11pt]{beamer}
\usetheme{Madrid}
\usefonttheme{serif}
\definecolor{customcolor}{RGB}{52,25,100} 
\setbeamercolor{structure}{fg=customcolor}

\usepackage[utf8]{inputenc}
\usepackage[T1]{fontenc}
\usepackage{amsmath}
\usepackage{MnSymbol}
\usepackage{amsfonts}
\usepackage{amssymb}
\usepackage{graphicx}
\usepackage{algorithm}
\usepackage{algpseudocode}
\usepackage{geometry}

\DeclareMathOperator{\sen}{sen}
\DeclareMathOperator{\tg}{tg}

\setbeamertemplate{caption}[numbered]

\author[Atishaya Maharjan]{Atishaya Maharjan}
\title{Geometric Spanning Trees that Minimizes the Wiener Index}
\newcommand{\email}{maharjaa@umanitoba.ca}
\setbeamertemplate{navigation symbols}{} 
\institute[]{University of Manitoba\par Geometric, Approximation, and Distributed Algorithms (GADA) lab} 
\date{\today} 

\bibliographystyle{apalike}

\begin{document}

\begin{frame}
	\titlepage
\end{frame}

\begin{frame}{This week}
	This week, I aimed to do the following:
	\begin{itemize}
		\item Read through the proof on the constant time approximation algorithm for the hamiltonian path and see if it already solves my problem.
		      \pause
		\item Run through and understand the NP-Hardness proof given for the Euclidean Minimum Spanning Tree in general positions with a wiener index and weight bound.
		      \pause
		\item Think about open problems relating to the Wiener index.
	\end{itemize}
\end{frame}

\begin{frame}{Constant time approximation algorithm for hamiltonian path minimizing the wiener index}
	\begin{definition}[k-spider]
		A \textbf{k-spider} is a tree with all vertices except the \textit{center} having degrees 1 or 2.
	\end{definition}

	\begin{definition}[k-repairman problem]
		Given:
		\begin{itemize}
			\item A metric space $(V, d)$ where $V$ is the set of points and $d$ is the distance function satisfying the triangle inequality.
			\item A positive integer $k$, representing the number of repairmen.
			\item The goal is to serve all clients using $k$ paths (one for each repairman), where each path starts from a designated origin or root.
		\end{itemize}
	\end{definition}
\end{frame}

\begin{frame}{Constant time approximation algorithm for hamiltonian path minimizing the wiener index}
	\begin{definition}[k-repairman problem (Contd)]
		Let $P_1, P_2, \dots, P_k$ be the paths assigned to the $k$ repairmen, and each client is visited exactly once (by one path), then:
		$$\min \dfrac{1}{|V|}\displaystyle\sum_{v\in V}t(v)$$
		where $t(v)$ is the arrival time of the node $v$ on the path it appears in.
	\end{definition}

	\begin{theorem}[Theorem 2.5 \cite{proceeding:approx_alg_for_min_avg_distortion}]
		Givem a $\gamma$-approximation for the minimum k-repairmen problem on a metric $d$, we can obtain a $2\gamma$-approximation for embedding the metric $d$ into a $k$-spider in a non-contracting fashion to minimize the average distortion.
	\end{theorem}
\end{frame}

\begin{frame}{Constant time approximation algorithm for hamiltonian path minimizing the wiener index}
	\begin{itemize}
		\item \cite{proceeding:approx_alg_for_min_avg_distortion} showed that embedding a metric into a line (So the sum of pairwise distances in the line - The Wiener index of the path is minimized) can be approximated via the k-TRP.
		\item Moreover, \cite{article:poly_time_approx_for_trp} gave a PTAS $(1 + \varepsilon)$-approximation for the Euclidean k-TRP (and for planar graphs). Plugging $\gamma = 1 + \varepsilon$ provides a $2 + \mathcal{O}(\varepsilon)$ approximation factor.
	\end{itemize}
\end{frame}

\begin{frame}{Weak vs Strong NP-Hardness reductions}
	We begin with two related problems:
	\begin{definition}[2-PARTITION]
		Given $n$ integers $A = \{a_1, a_2, \dots, a_n\}$, partition them into two groups $A_1$ and $A_2$, where $A_1 \cupdot A_2 = A$ such that $\displaystyle\sum_{i \in A_1}i = \displaystyle\sum_{j \in A_2}j = \dfrac{1}{2}\displaystyle\sum_{a \in A}a = t$.
	\end{definition}

	\begin{definition}[3-PARTITION]
		Given $n$ integers $A = \{a_1, a_2, \dots, a_n\}$ and no target sum, parition them into $\dfrac{n}{3}$ sets $A_1 \cupdot A_2 \cupdot \dots \cupdot A_{\frac{n}{3}} = A$ of equal sum: $\displaystyle\sum A_i = \dfrac{\displaystyle\sum A}{n/3} = t$. We will assume $3\;|\;n$; otherwise the answer is no.
	\end{definition}
\end{frame}

\begin{frame}{Weak vs Strong NP-Hardness reductions}
	\begin{itemize}
		\item Weakly NP-hard: NP-hard ``in the intended way''. Usually, we think of NP-hardness in terms of how hard your problem is, in terms of the encoding size $(n)$ of your input. Thus, we allow numbers to have exponential value because even when a value is exponential, the encoding of the value is polynomial (Read in binary: $\log_2(\text{number} = 2^{n^c}) = n^c$, which is polynomial). Polynomial number of bits is in some sense a reasonable encoding.

		\item Strongly NP-hard: NP-hard even when problem is restricted to numbers of polynomial value (With respect to n, the number of numbers).
	\end{itemize}

	\vspace{0.5em}
	Key: \textbf{2-PARTITION} is Weakly NP-hard while \textbf{3-PARTITION} is Strongly NP-hard.
\end{frame}

\begin{frame}{NP-Hardness reduction}
	\begin{definition}[Euclidean Wiener Index Tree Problem (EWITP)]
		Given a set $P$ of points in the Euclidean plane, a cost $W$, and a budget $B$, decide whether there exists a spanning tree $T$ of $P$, such tghat $W(T) \leq W$ and $wt(T) \leq B$.
	\end{definition}

	\begin{theorem}[Theorem 3 \cite{article:geometric_spanning_trees_minimizing_wiener_index}]
		The EWITP is weakly NP-hard since they reduced it from \textbf{2-PARTITION}.
	\end{theorem}
\end{frame}

\begin{frame}{Open problems from the geometric spanning tree paper}
	\begin{itemize}
		\item Improve the running time of $\mathcal{O}(n^4)$ algorithm that constructs a spanning tree of minimum Wiener index when the points of $P$ are in convex position.

		\item Is the problem of computing a spanning tree of minimum Wiener index for a set $P$ of $n$ points in general position NP-hard? The paper considered the variant where the total weight of the tree is bounded. If the total weight of the tree is unbounded, is it still NP-hard?
	\end{itemize}
\end{frame}

\begin{frame}{Next week tasks!}
	\begin{itemize}
		\item Read and understand the hardness reduction and understand why the bounding the weight of the tree is necessary.
		\item Think about either an efficient algorithm or a hardness reduction for the spanning tree minimizing the Wiener Index in general position.
	\end{itemize}
\end{frame}

\begin{frame}[allowframebreaks]{References}
	\bibliography{references}
\end{frame}

\begin{frame}
	\begin{center}
		The end.

		\email
	\end{center}
\end{frame}

\end{document}
